% !TEX root = ../../main.tex
% C7.2 — Dữ liệu thử nghiệm và camera giả lập. Số liệu: report-data-guide.md mục 2 (dữ liệu demo 1.488 / 1.523 track) và mục 3 C7.2 (Hiện hành, R6).
% Không lấy số từ database đang chạy. WILDTRACK: chavdarova2018wildtrack (7 camera, 1920x1080, ~60 fps).
\section{Dữ liệu thử nghiệm và camera giả lập}
\label{sec:7.2}

\paragraph{Bộ dữ liệu.} Đề tài sử dụng bộ dữ liệu WILDTRACK~\cite{chavdarova2018wildtrack}, gồm video của bảy camera tĩnh độ phân giải 1920$\times$1080 cùng quan sát một khu vực đông người đi bộ ngoài trời, với vùng nhìn chồng lấn nhau, kèm nhãn vị trí người trên một tập khung hình. Mỗi video được xem là một camera, và cả bảy camera được gán vào cùng một khu vực giám sát.

\paragraph{Camera giả lập.} Bảy video được phát thành bảy luồng RTSP qua MediaMTX như Mục~\ref{sec:3.7.3}. Tuy nhiên, trên máy triển khai, quá trình phân tích chậm hơn thời gian thực khoảng bảy lần, nên máy phát luồng bỏ bớt khung hình và phần lớn nội dung của luồng không được phân tích. Vì vậy, luồng RTSP được dùng để trình diễn và kiểm chứng việc tiếp nhận camera (Mục~\ref{sec:8.6}), còn dữ liệu trình diễn được lập chỉ mục từ chính các tệp video qua chức năng tải video lên, đi qua cùng quy trình xử lý.

\paragraph{Chuẩn bị dữ liệu trình diễn.} Việc lập chỉ mục được thực hiện bằng một công cụ gọi API của hệ thống như Quản trị viên: công cụ cắt đoạn video cần xử lý, tải lên cho từng camera với cùng một mốc thời gian gốc cho cả bảy camera, rồi theo dõi tới khi công việc kết thúc. Dữ liệu trình diễn phủ 120 giây đầu của mỗi video: đoạn 0-60 giây tạo 736 track, đoạn 60-120 giây được xử lý bằng 7 công việc trong 55 phút, tạo thêm 752 track, không có lỗi. Tổng cộng có 1.488 track ở trạng thái sẵn sàng, phân bố như Bảng~\ref{tab:demo-tracks}. Cùng với 35 track từ một phiên xử lý luồng RTSP, dữ liệu trình diễn có 1.523 track và một vụ việc mẫu gồm 3 kết quả.

\begin{table}[H]
\centering
\caption{Số track của dữ liệu trình diễn theo camera (120 giây đầu mỗi video)}
\label{tab:demo-tracks}
\begin{tabularx}{\textwidth}{|L|c|c|c|c|c|c|c|c|}
\hline
\textbf{Camera} & C1 & C2 & C3 & C4 & C5 & C6 & C7 & \textbf{Tổng} \\
\hline
\textbf{Số track} & 208 & 254 & 308 & 188 & 182 & 159 & 189 & 1.488 \\
\hline
\end{tabularx}
\end{table}
